% !TeX root = ../main.tex
\begin{frame}{03 创新思路——管状结构扩张响应}
  \begin{SlideCanvas}
    \BodyAt{55}{84}{850}{课题组通过改变目标泊松比，在相近轴向拉力响应下获得不同的径向扩张幅度。}
    \PointAt{55}{156}{245}{代表构型}
    \PointAt{343}{156}{260}{轴向拉伸力}
    \PointAt{639}{156}{270}{扩张半径}
    \PicAt{55}{202}{246}{239}{assets/group14_geometry.pdf}
    \PicAt{342}{202}{263}{239}{assets/group14_force.pdf}
    \PicAt{638}{202}{264}{239}{assets/group14_radius.pdf}
    \DetailAt{55}{449}{850}{目标泊松比变化提供可控的扩张幅度，与上一组算例构成互补的联合编程结果。}
    \Sources{本课题组成果：\RefGroup，图14(a)(e)(f)。目标负泊松比在$-0.8$至$-0.4$范围内变化。}
  \end{SlideCanvas}
  \BICConclusionBand{载荷近似保持时，\BICBandEmph{径向扩张幅度可以调整}}
  \note{本页仍然是轴向拉伸算例，不能把其力响应直接当作径向力—直径曲线。图14对应同一非线性响应目标与不同负泊松比目标。相近载荷响应和不同扩张曲线构成与上一页互补的证据，支持联合设计具有一定的协调能力。原文同样指出目标之间可能存在取舍。}
\end{frame}
